\section{Step 4 Return the full gcd pattern}\label{rsa-s04}

Move the final first value into \texttt{v0} and set \texttt{v1} to zero. In this interface \texttt{v0} is numerical data, so \texttt{v1} distinguishes the successful gcd return from the starter's unfinished return. Inspect the returned bits as an unsigned 64-bit value; do not narrow them to a signed word.

Replace only the body of region \texttt{gcd\_return}, retaining its BEGIN and END comments.

\Needspace{10\baselineskip}
\MARSCodeLabel{MIPS assembly}
\begin{minted}[samepage=true]{mars-idea}
gcd_done:
    move $v0, $t0
    move $v1, $zero
    jr $ra
    nop
\end{minted}

Use recipe \texttt{s04-1} with the cumulative source through this edit. Stop at \texttt{function\_return} before continuing to the temporary or normal return. With the supplied caller, use \texttt{gcd64\_return} as that return breakpoint.

\subsection{Expected observation}
The function returns gcd 1 and \texttt{v1} zero, preserving the caller's saved registers.

\Needspace{9\baselineskip}
\subsection{Check your understanding}
Which value indicates successful gcd completion?

\begin{itemize}
\item \textbf{A} The instruction count.
\item \textbf{B} Zero in \texttt{v1}.
\item \textbf{C} Zero in the gcd result for every input.
\end{itemize}
